% 张三慧《大学物理学》（第4版）力学、热学，第3章动量与角动量
% 章末习题3.1--3.20：原书PDF物理页116--117，印刷页100--101。
% 3.5和3.14题跨页，已连续检查；3.16、3.18依赖原书插图。
% 本文件是候选LaTeX片段，版权许可及独立学术审核未完成。
\section*{第3章\quad 动量与角动量}
\subsection*{章末习题（3.1--3.20，后续待录）}

\paragraph{习题 3.1}
一小球在弹簧的作用下振动（原书图3.34），弹力
\(F=-kx\)，而位移 \(x=A\cos\omega t\)，
其中 \(k,A,\omega\) 都是常量。
求在 \(t=0\) 到 \(t=\pi/(2\omega)\) 的时间间隔内
弹力施于小球的冲量。

\paragraph{习题 3.2}
一个质量 \(m=50\,\mathrm g\)，以速率
\(v=20\,\mathrm{m/s}\) 作匀速圆周运动的小球，
在 \(1/4\) 周期内向心力加给它的冲量是多大？

\paragraph{习题 3.3}
美国丹佛市每年举办一次“水世界肚皮水花比赛”，
2007年6月21日比赛参加者 Hoffman（冠军）
跳水的姿态见原书图3.35。设他的质量为
\(150\,\mathrm{kg}\)，跳起后离水面最大高度是
\(5.0\,\mathrm m\)，碰到水面 \(0.30\,\mathrm s\)
后开始缓慢下沉，求他碰水的力多大？
\textit{理想化：}碰水时平均阻力近似恒定，
并在 \(0.30\,\mathrm s\) 中使竖直向下速度减至近零。

\paragraph{习题 3.4}
自动步枪连续发射每分钟射出 \(120\) 发子弹，
每发子弹的质量为 \(m=7.90\,\mathrm g\)，
出口速率为 \(735\,\mathrm{m/s}\)。
求射击时（以分钟计）枪托对肩部的平均压力。

\paragraph{习题 3.5}
2007年2月28日凌晨2时，由乌鲁木齐开出的
5807次旅客列车在吐鲁番境内突然遭遇13级特大飓风袭击，
11节车厢脱轨（原书图3.36），造成了3死34伤的惨祸。
13级飓风风速按 \(137\,\mathrm{km/h}\)（\(38\,\mathrm{m/s}\)）计，
空气密度为 \(1.29\,\mathrm{kg/m^3}\)，
车厢长 \(22.5\,\mathrm m\)，高 \(2.62\,\mathrm m\)。
设大风垂直吹向车厢侧面，碰到车厢后就停下来。
求飓风对一节车厢的水平推力多大？

\paragraph{习题 3.6}
水管有一段弯曲成 \(90^\circ\)。
已知管中水的流量为
\(3\times10^3\,\mathrm{kg/s}\)，流速为
\(10\,\mathrm{m/s}\)。
求水流对这弯管的压力的大小和方向。

\paragraph{习题 3.7}
桌面上堆放一串柔软的长链，今拉住长链的一端竖直向上
以恒定速度 \(v_0\) 上提。
试证明：当提起的长度为 \(l\) 时，
所用的向上拉力
\[
 F=\rho_l lg+\rho_l v_0^2,
\]
其中 \(\rho_l\) 为长链单位长度的质量。

\paragraph{习题 3.8}
手提住一柔软长链的上端，
使其下端刚与桌面接触，然后松手使链自由下落。
试证明下落过程中，桌面受的压力等于
已落在桌面上的链的重量的3倍。

\paragraph{习题 3.9}
一个原来静止的原子核，放射性衰变时放出一个动量为
\(p_1=9.22\times10^{-21}\,\mathrm{kg\,m/s}\) 的电子，
同时还在垂直于此电子运动的方向上放出一个动量为
\(p_2=5.33\times10^{-21}\,\mathrm{kg\,m/s}\) 的中微子。
求衰变后原子核的动量的大小和方向。

\paragraph{习题 3.10}
运载火箭的最后一级以
\(v_0=7600\,\mathrm{m/s}\) 的速率飞行。
这一级中一个质量为 \(m_1=290.0\,\mathrm{kg}\)
的火箭壳和一个质量为 \(m_2=150.0\,\mathrm{kg}\)
的仪器舱抱在一起。当抛松开后，
二者间的压缩弹簧使二者分离，
这时二者的相对速率为
\(u=910.0\,\mathrm{m/s}\)。
设所有速度都在同一直线上，
求两部分分开后各自的速度。

\paragraph{习题 3.11}
两辆质量相同的汽车在十字路口垂直相撞，
撞后二者扣在一起沿直线滑动了
\(s=25\,\mathrm m\) 才停下来。
设滑动时地面与车轮之间的滑动摩擦系数为
\(\mu_k=0.80\)。
撞后两个司机都声明在撞车前自己的车速未超限
（\(14\,\mathrm{m/s}\)），他们的话都可信吗？

\paragraph{习题 3.12}
一空间探测器质量为 \(6090\,\mathrm{kg}\)，
正相对于太阳以 \(105\,\mathrm{m/s}\) 的速率向木星运动。
当它的火箭发动机相对于它以
\(253\,\mathrm{m/s}\) 的速率向后喷出
\(80.0\,\mathrm{kg}\) 废气后，
它对太阳的速率变为多少？
\textit{物理页117图像校读警示：}题中质量应与原页的
\textbf{6090 kg}核对；本稿保留原扫描读数。

\paragraph{习题 3.13}
在太空静止的一单级火箭，
点火后，其质量的减少与初质量之比为多大时，
它喷出的废气将是静止的？

\paragraph{习题 3.14}
一质量为 \(2.72\times10^6\,\mathrm{kg}\) 的火箭竖直离地面发射，
燃料燃烧速率为 \(1.29\times10^3\,\mathrm{kg/s}\)。
\begin{enumerate}
\item 它喷出的气体相对于火箭体的速率是多大，
才能使火箭刚刚离开地面？
\item 它以恒定相对速率 \(5.50\times10^4\,\mathrm{m/s}\)
 喷出废气，全部燃烧时间为 \(155\,\mathrm s\)。
 它的最大上升速率多大？
\item 在(2)的情形下，当燃料刚燃烧完时，
火箭体离地面多高？
\end{enumerate}
\textit{使用原题近似：}重力加速度视为常量，
忽略空气阻力。

\paragraph{习题 3.15}
一架喷气式飞机以 \(210\,\mathrm{m/s}\) 的速度飞行，
它的发动机每秒钟吸入 \(75\,\mathrm{kg}\) 空气，
在体内与 \(3.0\,\mathrm{kg}\) 燃料燃烧后
以相对于飞机 \(490\,\mathrm{m/s}\) 的速度向后喷出。
求发动机对飞机的推力。

\paragraph{习题 3.16}
水分子的结构如原书图3.37所示。
两个氢原子与氧原子的中心距离都是
\(0.0958\,\mathrm{nm}\)，它们与氧原子中心的连线的夹角为
\(105^\circ\)。求水分子的质心。

\paragraph{习题 3.17}
求半圆形均匀薄板的质心。

\paragraph{习题 3.18}
有一正立方体铜块，边长为 \(a\)。
今在其下半部中央挖去一截面半径为 \(a/4\) 的圆柱形洞
（原书图3.38）。求剩余铜块的质心位置。
\textit{图示依赖：}圆柱孔沿立方体深度方向贯穿，
孔轴位于截面中央偏下方，准确轴线位置须和图3.38复核。

\paragraph{习题 3.19}
在楼顶释放一质量 \(m_1=20\,\mathrm g\) 的石子后，
\(1\,\mathrm s\) 末又自同一点释放另一质量
\(m_2=50\,\mathrm g\) 的石子。
求在前者释放后 \(t\ (>1)\,\mathrm s\)，
这两个石子系统的质心的速度和加速度。

\paragraph{习题 3.20}
哈雷彗星绕太阳运动的轨道是一个椭圆。
它离太阳最近的距离是
\(r_1=8.75\times10^{10}\,\mathrm m\)，
此时它的速率是
\(v_1=5.46\times10^4\,\mathrm{m/s}\)。
它离太阳最远时的速率是
\(v_2=9.08\times10^2\,\mathrm{m/s}\)，
这时它离太阳的距离 \(r_2\) 是多少？
